% =====================================================================
%  TITRE III : DE L'ACTION OPÉRATIONNELLE
% =====================================================================
\titrepartie{Titre III}{De l'action opérationnelle}

% ---------------------------------------------------------------------
\article{Discipline radio et communications}

\chapeau{La radio est l'outil de commandement du département. Sa saturation coûte des secondes qui peuvent coûter une vie. Sa discipline n'est pas une question de forme.}

\cl[radio-principe]{La fréquence de service est réservée au trafic opérationnel. Aucune conversation étrangère au service n'y est admise. Toute discussion appelée à durer bascule sur une fréquence secondaire, sur demande de l'une des unités concernées.}

\cl[format]{Toute transmission respecte le format suivant : indicatif de l'émetteur, destinataire, message. L'émetteur attend l'accusé de réception avant de poursuivre. Une transmission porte une information : les messages composites sont scindés.}

\cl{L'agent formule mentalement son message avant d'ouvrir le micro. Les hésitations, reprises et corrections en direct saturent la fréquence.}

\cl[espacement]{Un intervalle de trois à cinq secondes est observé entre deux transmissions afin de permettre à une unité en urgence de percer.}

\cl[priorite]{Sur annonce de trafic prioritaire par une unité ou par le Watch Commander, toutes les unités non concernées observent le silence radio jusqu'à la levée de la priorité.}

\cl{La prise de service, la fin de service, l'indisponibilité temporaire, l'arrivée sur les lieux et la fin d'intervention font l'objet d'une annonce systématique. Une unité dont la position n'est pas connue du Watch Commander est une unité perdue.}

\cl[code99]{L'annonce de \code{Code 99} est réservée à la situation d'un agent en danger immédiat. Elle déclenche la réponse de l'ensemble des unités disponibles en Code 3 et le silence radio général. Elle ne s'emploie ni pour un accident de la route, ni pour une interpellation difficile, ni pour un renfort de confort.}

\begin{interdit}[Sur la fréquence de service]
Sont interdits : l'usage de la radio pour insulter, plaisanter ou commenter une intervention ; l'emploi d'un code que l'agent ne maîtrise pas ; la diffusion d'informations nominatives non nécessaires ; le maintien du micro ouvert. L'usage abusif de la radio est une faute disciplinaire.
\end{interdit}

\begin{rappel}[Discipline de fréquence en intervention]
Sur une intervention en cours, seule l'unité primaire rend compte. Les autres unités s'annoncent une fois, se placent, et se taisent. Le Operation Lead est le seul interlocuteur du Watch Commander.
\end{rappel}

% ---------------------------------------------------------------------
\article{Avis de recherche}

\cl[bolo]{Un avis de recherche est diffusé lorsqu'une personne, un véhicule ou un bien doit être identifié ou localisé par l'ensemble des unités. Il est diffusé par l'unité qui en dispose et rediffusé, si nécessaire, par le Watch Commander.}

\sstitre{Contenu obligatoire}

\vspace{2pt}
\noindent
\begin{tabular}{P{30mm} P{118mm}}
\headrow
\thead{Objet} & \thead{Éléments à transmettre impérativement} \\
\tkey{Personne} & \tcell{Sexe, couleur de peau, description du haut (type et couleur), description du bas (type et couleur), au moins un signe distinctif, direction de fuite si connue, et \textbf{motif de la recherche}.} \\
\altrow
\tkey{Véhicule} & \tcell{Type, couleur, marque et modèle si connus, au moins trois caractères de l'immatriculation, nombre d'occupants, direction de fuite, et \textbf{motif de la recherche}.} \\
\tkey{Bien} & \tcell{Description précise de l'objet, éléments d'identification (numéro de série, marquage), et \textbf{motif de la recherche}.} \\
\bottomrule
\end{tabular}

\cl{L'unité émettrice joint systématiquement son indicatif afin que les unités disposant d'un élément sachent qui contacter.}

\cl[bolo-levee]{Tout avis de recherche est levé dès sa résolution, par annonce de l'unité émettrice. Le maintien d'un avis résolu sur la fréquence est une faute de procédure : il mobilise inutilement des unités et expose des citoyens à un contrôle injustifié.}

\cl{Un avis de recherche ne constitue ni un mandat, ni une autorisation d'interpellation. Il fonde une suspicion raisonnable permettant le contrôle ; l'interpellation suppose la vérification des éléments sur place.}

% ---------------------------------------------------------------------
\article{Pouvoir d'appréciation, suspicion raisonnable et cause probable}

\chapeau{Toute intervention repose sur un fondement juridique. Un agent qui ne peut pas énoncer le sien n'a pas de raison d'agir.}

\cl[suspicion]{La \stress{suspicion raisonnable} est constituée de faits précis et exprimables, et des déductions rationnelles qui peuvent en être tirées. Elle autorise le contrôle d'une personne ou d'un véhicule et la palpation de sécurité. Elle ne peut jamais reposer sur la seule apparence, le seul lieu de présence ou le seul antécédent d'une personne.}

\cl[cause-probable]{La \stress{cause probable} est constituée de faits qui conduiraient une personne raisonnable à croire qu'une infraction a été commise et que la personne visée en est l'auteur. Elle est requise pour procéder à une arrestation, à une fouille approfondie ou à une saisie.}

\begin{note}[Distinguer les deux notions]
Un véhicule dont le type et la couleur correspondent à celui signalé sur un fait commis deux minutes plus tôt à deux rues de là fonde une suspicion raisonnable : on contrôle. Un véhicule de couleur approchante, à cinq kilomètres et vingt minutes de l'événement, n'en fonde aucune. La suspicion permet de vérifier ; la cause probable permet de contraindre.
\end{note}

\cl[discretion]{L'agent dispose d'un pouvoir d'appréciation. Pour une infraction non violente et sans victime constituée, il peut délivrer un avertissement ou une contravention en lieu et place d'une interpellation. Ce choix est apprécié au regard du comportement de la personne, du risque créé et de l'intérêt du service.}

\cl[desescalade]{La désescalade précède l'escalade. Lorsque la situation le permet, l'agent parle, ralentit le rythme de l'échange, crée de la distance, met un obstacle entre lui et le sujet, et demande du renfort plutôt que d'engager seul. Le temps est presque toujours un allié.}

\cl{L'agent conserve la maîtrise du ton en toute circonstance. Il ne s'aligne jamais sur le registre de son interlocuteur, ne répond pas à la provocation et ne se laisse pas déborder. La provocation d'un civil se constate, se documente et se rapporte : elle ne s'entretient pas.}

\begin{interdit}[Interdictions permanentes]
Sont interdits : le contrôle d'une personne sans suspicion raisonnable énonçable ; la surveillance prolongée d'un civil dans l'attente d'une infraction ; le cumul artificiel de charges destiné à alourdir une procédure ; et toute décision fondée en tout ou partie sur l'apparence, l'origine, le sexe ou les croyances d'une personne.
\end{interdit}

% ---------------------------------------------------------------------
\article{Usage de la force}

\cl[force-habilitation]{Tout agent ayant achevé sa formation initiale est habilité à faire usage de la force dans les conditions et les limites fixées au présent article.}

\cl[proportionnalite]{L'usage de la force est mesuré, proportionné et gradué. L'agent recourt au niveau de contrainte immédiatement supérieur à celui opposé par le sujet, dans la limite du raisonnable, et redescend d'un niveau dès que la résistance cesse. Est raisonnable la force qu'un agent raisonnable jugerait nécessaire dans des circonstances identiques.}

\sstitre{Échelle graduée de la force}

\vspace{2pt}
\noindent
\begin{tabular}{C{14mm} P{40mm} P{88mm}}
\headrow
\thead{Niveau} & \thead{Moyen} & \thead{Emploi} \\
\tcell{\lspdnum{1}} & \tkey{Force de présence} & \tcell{La présence identifiée d'un agent suffit dans la majorité des situations. S'identifier comme policier est le premier acte de contrainte.} \\
\altrow
\tcell{\lspdnum{2}} & \tkey{Communication} & \tcell{Sommations, injonctions claires et répétées, négociation. Beaucoup de situations se désamorcent par le choix des mots.} \\
\tcell{\lspdnum{3}} & \tkey{Contrainte à mains nues} & \tcell{Techniques de contrôle et de maîtrise permettant de maîtriser un sujet récalcitrant sans recours à une arme.} \\
\altrow
\tcell{\lspdnum{4}} & \tkey{Force non létale} & \tcell{Taser, matraque, lanceur à balles de défense, projection de gaz. Résistance active ou agression physique.} \\
\tcell{\lspdnum{5}} & \tkey{Force létale} & \tcell{Menace de mort ou de blessure grave, imminente et actuelle. Dernier recours absolu.} \\
\bottomrule
\end{tabular}

\sstitre{Force létale}

\cl[letale]{La force létale ne peut être employée contre une personne que dans les cas suivants, et dans ces cas seulement :}
\begin{itemize}
  \item pour empêcher cette personne de tuer ou de blesser gravement un tiers ou un agent ;
  \item lorsqu'elle constitue une menace mortelle immédiate et actuelle pour la sécurité des citoyens ou des agents.
\end{itemize}

\cl[sommations]{L'agent adresse \textbf{trois sommations claires et distinctes} avant tout tir et avant tout emploi du taser. Il annonce son identité et son intention. Il n'est dispensé de cette obligation que lorsqu'un impératif contraire, tenant à la soudaineté de l'agression ou au danger de mort immédiat pour un tiers, rend l'avertissement matériellement impossible.}

\cl[identification]{L'agent identifie sa cible et son arrière-plan avant tout tir. Il ne tire pas lorsque la visibilité est obstruée, lorsqu'un otage ou un tiers se trouve dans l'axe, ou lorsque l'arrière-plan n'est pas maîtrisé. En cas de doute sur l'identification, l'agent ne tire pas.}

\cl{L'agent annonce clairement sa présence à ses collègues lorsqu'il apparaît dans leur champ de tir. Le tir fratricide est une faute d'identification, jamais une fatalité.}

\cl[fin-menace]{L'autorisation d'employer la force létale n'est jamais permanente : elle prend fin avec la menace. Dès que le sujet cesse de constituer un danger (désarmement, immobilisation, reddition, incapacité), le tir cesse et l'interpellation reprend son cours normal.}

\cl[rapport-force]{Tout usage de la force létale donne lieu à un rapport et à un dossier d'incident. Une enquête interne peut être ouverte à la demande d'un superviseur, sous accord du Command Staff. Cette décision fait l'objet d'une note de service précisant les éléments qui la justifient.}

\cl{Tout coup de feu tiré par un agent, qu'il ait ou non atteint sa cible, entraîne automatiquement : la sécurisation et la préservation de la scène, la conservation immédiate des enregistrements de caméra-piéton, la saisine du Detective Bureau et celle de la Division des Affaires Internes, ainsi que le retrait administratif temporaire de l'agent concerné. Ce retrait est une mesure conservatoire et ne préjuge d'aucune faute.}

\begin{rappel}[Repère de distance]
Face à une arme blanche, un sujet déterminé parcourt une dizaine de mètres avant qu'un agent n'ait pu dégainer et tirer. Créer et maintenir de la distance n'est pas une reculade : c'est la seule manière de conserver une option non létale.
\end{rappel}

% ---------------------------------------------------------------------
\article{Armes intermédiaires et armement spécialisé}

\sstitre{Taser}

\cl[taser]{L'emploi du taser est précédé de \textbf{trois sommations claires et distinctes}, audibles et sans ambiguïté, sauf impératif contraire au sens de l'article relatif à l'usage de la force. Le taser vise à obtenir la soumission d'un sujet ; il n'est ni un moyen de punition, ni un moyen d'arrêter une fuite véhiculaire.}

\cl[taser-interdit]{L'emploi du taser est interdit dans les situations suivantes :}
\begin{itemize}
  \item lorsqu'il apparaît que la décharge ou la chute qui s'ensuivra blessera gravement l'individu ;
  \item lorsque le sujet est déjà maîtrisé, menotté ou sous contrôle ;
  \item lorsque l'individu conduit un véhicule motorisé ;
  \item lorsque le sujet se trouve en hauteur, sur un toit, une corniche, une échelle, un escalier ou une pente ;
  \item lorsque le sujet se trouve dans l'eau ;
  \item lorsque des matières inflammables sont présentes à proximité.
\end{itemize}

\sstitre{Lanceur à balles de défense}

\cl[bean-bag]{L'emploi du lanceur requiert l'accord du Operation Lead ou d'un superviseur, sauf danger de mort immédiat. Il est précédé de trois sommations, annoncées aux autres agents présents afin d'éviter toute confusion avec un tir d'arme à feu.}

\cl{Le tir vise les cuisses et le bas de l'abdomen. Sont proscrits les tirs à la tête, au cou, au thorax et à l'aine. Une distance minimale de tir est respectée, sauf danger de mort immédiat. Le tir depuis une position surélevée ou vers un sujet situé en hauteur est interdit.}

\sstitre{Bouclier balistique}

\cl[bouclier]{Le déploiement du bouclier balistique requiert l'accord d'un superviseur. Son porteur ouvre la progression et entre le premier. Il n'est armé que d'une arme de poing.}

\sstitre{Fusil de précision}

\cl[precision]{Le fusil de précision est transporté dans le coffre du véhicule et son emploi est réservé aux agents qualifiés. Sa fonction principale est l'observation et le renseignement. Le tir létal n'est autorisé que sur ordre exprès du Operation Lead, dans une circonstance de danger de mort caractérisée.}

% ---------------------------------------------------------------------
\article{Contrôle routier}

\chapeau{Le contrôle routier est l'acte de police le plus fréquent et l'un des plus dangereux. Sa méthode n'est pas une préférence personnelle.}

\sstitre{Contrôle de routine}

\cl[controle-annonce]{Tout contrôle fait l'objet d'une annonce radio : position, description du véhicule, immatriculation, nombre d'occupants apparents et motif du contrôle. Cette annonce peut être faite avant l'arrêt ou, plus commodément, une fois le véhicule immobilisé et avant que l'agent ne descende de son véhicule de patrouille. Elle intervient dans tous les cas \textbf{avant tout contact avec les occupants}. Un contrôle non annoncé est un contrôle sans filet.}

\cl{L'agent déclenche ses signaux lumineux et guide le véhicule vers un emplacement sûr : bande d'arrêt d'urgence, aire de stationnement, voie secondaire dégagée. Il évite les carrefours, les voies rapides à fort trafic et les abords immédiats des établissements sensibles.}

\cl[positionnement]{Le véhicule de patrouille se positionne en retrait du véhicule contrôlé, moteur et feux allumés. Trois positions sont admises :}
\begin{itemize}
  \item \stress{en ligne} : protection minimale, réservée aux voies étroites ;
  \item \stress{décalé} : le véhicule empiète d'une demi-largeur sur la voie, protection intermédiaire ;
  \item \stress{en angle} : position de référence, offrant la meilleure protection au poste d'approche.
\end{itemize}

\cl{De nuit, l'agent emploie ses projecteurs de manière à éclairer l'habitacle sans s'exposer, et revêt un équipement à haute visibilité sur voie rapide.}

\sstitre{Mise en sécurité des occupants}

\cl[mise-en-securite]{Le département n'approche pas les véhicules contrôlés. \textbf{Ce sont les occupants qui viennent à l'agent, et non l'inverse.} Cette règle vaut pour l'ensemble des contrôles routiers, y compris les plus anodins : elle supprime l'angle mort de l'habitacle, qui est la principale cause d'agression d'agent lors d'un contrôle.}

\cl[sortie-occupants]{Depuis son véhicule de patrouille, l'agent donne au conducteur et à l'ensemble des passagers l'instruction de descendre du véhicule et de se placer sur le trottoir ou sur un accotement dégagé, à distance de la chaussée et de leur véhicule, mains visibles. Il indique clairement l'emplacement souhaité.}

\cl{L'agent ne descend de son véhicule qu'une fois l'ensemble des occupants sortis et positionnés à l'emplacement indiqué. Il conserve la maîtrise visuelle du véhicule contrôlé pendant toute la manœuvre.}

\cl{Lorsqu'un occupant refuse de descendre, l'agent ne force ni l'ouverture ni le contact : il maintient sa position, demande du renfort et bascule sur la procédure de contrôle à haut risque.}

\cl{Un occupant dont l'état physique ne permet pas la descente du véhicule est laissé à sa place ; l'agent en tient compte dans son approche et le signale à la radio.}

\sstitre{Contact}

\cl[contact]{L'agent se présente, annonce son service et \textbf{énonce le motif du contrôle} avant toute autre demande. Il sollicite ensuite les documents : permis de conduire, certificat d'immatriculation et attestation d'assurance.}

\cl{L'agent maintient les occupants dans son champ de vision et se place de manière à ne jamais avoir un occupant dans son dos.}

\cl{Le retour au véhicule de patrouille s'effectue de profil, sans jamais tourner le dos aux occupants ni au véhicule contrôlé.}

\cl{À l'issue du contrôle, l'agent restitue les documents, notifie la suite donnée (avertissement, contravention ou interpellation), autorise les occupants à regagner leur véhicule et annonce la fin de l'intervention à la radio.}

\sstitre{Contrôle à haut risque}

\cl[felony]{Le contrôle à haut risque s'impose notamment lorsque : le véhicule est signalé volé, un occupant est signalé armé, le contrôle fait suite à une poursuite, ou un occupant fait l'objet d'un mandat.}

\cl[felony-renfort]{Aucun contrôle à haut risque n'est engagé sans renfort suffisant. L'unité attend au minimum une seconde unité, sauf danger immédiat pour un tiers. L'arrêt peut être différé afin d'être conduit dans un lieu dégagé, à l'écart des axes passants et des zones piétonnes.}

\cl[felony-procedure]{Aucun agent n'approche le véhicule. Un seul agent donne les ordres, au moyen du haut-parleur ; les autres assurent la couverture depuis un abri. Les ordres sont donnés au conducteur dans l'ordre suivant :}
\begin{enumerate}
  \item couper le moteur et poser les clés sur le toit ou les jeter hors du véhicule ;
  \item ouvrir la portière depuis l'extérieur et sortir, mains sur la tête, dos aux agents ;
  \item reculer lentement vers les agents jusqu'à l'ordre d'arrêt ;
  \item se mettre à genoux, mains en évidence.
\end{enumerate}

\cl{Les passagers reçoivent l'ordre de demeurer immobiles dans le véhicule, puis sont extraits un par un selon la même séquence. Deux agents désignés procèdent au menottage, à la palpation et à l'évacuation vers l'arrière du dispositif ; les autres maintiennent la couverture sur le véhicule.}

\cl{Le véhicule n'est approché qu'une fois vidé de ses occupants et fait l'objet d'un contrôle visuel avant toute fouille. Les éléments susceptibles de valeur probante ne sont ni déplacés ni manipulés avant leur relevé.}

\begin{interdit}[Erreurs constitutives d'une faute]
Engager seul un contrôle à haut risque sans danger immédiat pour un tiers ; approcher un véhicule occupé ; donner des ordres à plusieurs voix simultanément ; extraire plusieurs occupants en même temps ; omettre l'annonce radio.
\end{interdit}

% ---------------------------------------------------------------------
\article{Interpellation, notification des droits et procédure d'arrestation}

\chapeau{Une interpellation irrégulière est une interpellation perdue. La procédure décrite ci-après est la seule admise.}

\sstitre{Notification des droits}

\cl[miranda]{Avant tout interrogatoire, et au plus tard dans les quinze minutes suivant l'interpellation, l'agent notifie ses droits à la personne interpellée. Il énonce au préalable, à voix haute, \textbf{la date, l'heure et le motif de l'arrestation}, puis lit la formule suivante :}

\begin{formule}
\centering\fontsize{10}{15.5}\selectfont
« Il est [heure], le [date]. Vous êtes en état d'arrestation pour [motif].\\[5pt]
Vous avez le droit de garder le silence.\\
Tout ce que vous direz pourra et sera retenu contre vous devant un tribunal.\\
Vous avez le droit à l'assistance d'un avocat.\\
Si vous n'avez pas les moyens d'en engager un, un avocat vous sera commis d'office.\\
Vous avez droit à de l'eau, à de la nourriture et à des soins médicaux.\\[5pt]
Comprenez-vous ces droits tels que je viens de vous les énoncer ?\\
Compte tenu de ces droits, acceptez-vous de vous entretenir avec moi\\ ou avec un autre agent des forces de l'ordre ? »
\end{formule}

\cl{L'agent s'assure que la personne a compris ses droits et consigne sa réponse. À défaut de réponse, le silence vaut refus de s'entretenir : l'interrogatoire ne peut être engagé.}

\cl[miranda-repetition]{Lorsque la personne ne répond pas ou déclare ne pas comprendre ses droits, la formule lui est lue \textbf{trois fois au total}. À l'issue de la troisième lecture, les droits sont réputés notifiés et acquis, et la procédure se poursuit. Le nombre de lectures est consigné au rapport d'arrestation.}

\sstitre{Déroulement de l'interpellation}

\begin{enumerate}
  \item \stress{Fondement} : l'agent dispose d'une cause probable ou d'un mandat valide. Il énonce à voix haute le motif de l'interpellation.
  \item \stress{Sommations} : trois sommations claires et distinctes précèdent tout recours à la force, sauf impératif contraire.
  \item \stress{Menottage} : mains dans le dos, vérification du serrage, annonce explicite de la mise en état d'arrestation.
  \item \stress{Palpation de sécurité} : systématique avant toute mise en véhicule, et renouvelée après tout passage en établissement de soins.
  \item \stress{Notification des droits} : dans les conditions prévues ci-dessus.
  \item \stress{Transport} : en Code 1, avec annonce radio du départ, de la destination et de l'arrivée.
  \item \stress{Inventaire} : inventaire contradictoire des effets saisis, remise d'un reçu à l'intéressé, dépôt sous scellés.
  \item \stress{Rapport} : rédaction du rapport d'arrestation avant la fin de service.
\end{enumerate}

\cl[palpation]{La palpation est effectuée, dans la mesure du possible, par un agent du même sexe que la personne contrôlée. À défaut, elle s'effectue exclusivement par-dessus les vêtements, sous caméra, et après annonce verbale enregistrée.}

\cl[detention-courte]{Une personne retenue aux fins de vérification sans qu'aucun élément à charge ne soit établi est remise en liberté dans un délai maximal de quinze minutes.}

\cl[charges]{Seules les charges prévues par le Code pénal peuvent être retenues. Aucune charge ne peut être ajoutée après le prononcé de la peine. Le cumul artificiel de charges destiné à alourdir la procédure est prohibé et constitue une faute disciplinaire.}

\begin{interdit}[Vices de procédure]
Entraînent la nullité de la procédure et engagent la responsabilité de l'agent : l'absence de notification des droits avant interrogatoire ; l'absence d'énoncé du motif de l'interpellation ; l'absence de palpation avant mise en véhicule ; le maintien d'effets illicites en possession du détenu ; l'absence de rapport d'arrestation.
\end{interdit}

% ---------------------------------------------------------------------
\article{Transport de détenus et escortes}

\cl[transport]{Tout transport de détenu s'effectue en \code{Code 1}. Le véhicule ne fait usage ni des signaux lumineux ni de la sirène et respecte intégralement le Code de la route. Un détenu n'est pas une urgence.}

\cl{Une palpation est effectuée avant toute mise en véhicule, y compris lorsqu'un autre agent y a déjà procédé. Aucune exception.}

\cl{Le départ et l'arrivée du transport sont annoncés à la radio, avec l'indication du point de départ, de la destination et du nombre de détenus transportés.}

\cl[surveillance-transport]{Un détenu n'est jamais laissé seul dans un véhicule, ni hors de la vue d'un agent, à aucun moment du transport. Lorsque la séparation est matériellement possible, deux mis en cause d'une même affaire ne voyagent pas ensemble.}

\cl{Tout détenu blessé est conduit sous escorte vers un établissement de soins avant sa présentation en cellule. Une palpation est renouvelée à la sortie de l'établissement.}

\sstitre{Escorte de convoi}

\cl[convoi]{Une escorte de convoi comprend au minimum trois unités : un véhicule de tête, le véhicule transportant les détenus ou les valeurs, et un véhicule de queue. Elle est commandée par un superviseur.}

\cl{L'itinéraire n'est pas diffusé sur la fréquence ouverte et est modifié dès qu'il est susceptible d'avoir été compromis.}

\cl[convoi-attaque]{En cas d'attaque du convoi, les véhicules d'escorte ne rompent jamais la protection du véhicule principal pour engager les assaillants. Le véhicule principal quitte la zone sans délai ; les renforts prennent en charge l'agression.}

\cl{Une fois le convoi parvenu à destination et les détenus remis, aucune libération ne peut plus intervenir du fait du département.}

% ---------------------------------------------------------------------
\article{Garde à vue et droits du détenu}

\cl[gav-duree]{La durée de la garde à vue est fixée à \textbf{une heure} pour un délit ordinaire. Elle peut être prolongée, sur autorisation d'un magistrat, jusqu'à vingt-quatre heures pour un crime, et au-delà pour les affaires relevant de la sécurité de l'État.}

\cl[gav-droits]{Dès son placement en garde à vue, et au plus tard dans les quinze minutes, le détenu est informé des droits suivants :}
\begin{itemize}
  \item le droit de garder le silence ;
  \item le droit d'être informé précisément des faits qui lui sont reprochés ;
  \item le droit à l'assistance d'un avocat ;
  \item le droit à un examen médical ;
  \item le droit à l'eau et à l'alimentation ;
  \item le droit de faire prévenir un proche.
\end{itemize}

\cl[avocat]{L'avocat sollicité dispose d'un délai raisonnable pour se manifester. Une fois présent, il bénéficie d'un entretien confidentiel avec son client. \textbf{Aucun agent n'est autorisé à assister à cet entretien, à l'écouter ou à l'enregistrer.}}

\cl{L'accès de l'avocat à son client ne peut être refusé que si sa présence physique compromet la sécurité des lieux. Ce refus est motivé par écrit et porté sans délai à la connaissance d'un superviseur.}

\cl[fouille-corps]{La fouille à corps est effectuée par un agent du même sexe que la personne fouillée. À défaut, elle s'effectue exclusivement par-dessus les vêtements et sous caméra. La fouille intégrale requiert l'accord d'un superviseur et est consignée.}

\cl[perquisition]{Toute perquisition requiert un mandat validé, sauf flagrance caractérisée. Toute saisie donne lieu à un procès-verbal détaillé et contradictoire.}

\cl{La fouille d'un téléphone, d'un ordinateur ou de tout support de données requiert le consentement écrit de son propriétaire ou un mandat de perquisition exprès.}

\cl[surveillance-cellule]{Un détenu n'est jamais laissé sans surveillance. Des rondes régulières sont effectuées et consignées. Tout incident survenu en cellule fait l'objet d'un rapport immédiat.}

\cl{Les mineurs et les personnes en situation de vulnérabilité manifeste ne peuvent être entendus hors la présence d'un responsable ou d'un tiers qualifié, qui est avisé sans délai.}

\begin{interdit}[Pratiques prohibées en garde à vue]
Prolonger une détention dans le seul but d'exercer une pression ; conduire un interrogatoire sans notification préalable des droits ; laisser un objet illicite en possession d'un détenu ; refuser ou différer sans motif l'accès à un avocat ou à un examen médical ; refuser de recevoir une doléance.
\end{interdit}

% ---------------------------------------------------------------------
\article{Caméra-piéton, preuves et scellés}

\cl[bodycam]{La caméra-piéton est activée pendant l'intégralité de toute interaction avec le public : contrôle, interpellation, palpation, fouille, usage de la force, audition. L'agent annonce à son interlocuteur que l'interaction est enregistrée.}

\begin{interdit}[Enregistrements]
Il est interdit d'interrompre, d'effacer, de modifier ou de dissimuler un enregistrement de caméra-piéton. Cette faute est lourde : elle entraîne la révocation et la saisine de la Division des Affaires Internes, indépendamment du contenu de l'enregistrement.
\end{interdit}

\cl{Tout usage de la force et tout coup de feu entraînent la conservation immédiate de la séquence correspondante et son versement au dossier de procédure.}

\cl[scelles]{Toute preuve saisie est photographiée sur place, référencée (numéro de scellé, agent saisissant, date, heure et lieu de la saisie) puis déposée en salle des scellés. Chaque manipulation ultérieure est tracée et signée.}

\cl[chaine-preuve]{Une preuve dont la chaîne de conservation est rompue est irrecevable. La rupture de la chaîne engage la responsabilité de l'agent qui en est à l'origine.}

\cl{Les enregistrements et les pièces de procédure ne sortent du département que sur réquisition régulière de l'autorité judiciaire ou sur demande d'un avocat transmise par la voie institutionnelle.}

% ---------------------------------------------------------------------
\article{Rapports et documentation}

\begin{solennel}
\centering\normalsize
{\lspddisp\fontsize{12}{15}\selectfont\color{lspdnavy}\upshape
Ce qui n'est pas écrit n'a pas eu lieu.}
\end{solennel}

\cl[rapports]{Les rapports constituent la mémoire du département et le fondement de toute procédure ultérieure. Leur rédaction n'est pas une formalité administrative : c'est un acte de police.}

\vspace{4pt}
\noindent
\begin{tabular}{P{54mm} P{56mm} P{36mm}}
\headrow
\thead{Rapport} & \thead{Fait générateur} & \thead{Délai} \\
\tcell{Rapport d'arrestation} & \tcell{Toute interpellation} & \tcell{Avant la fin de service} \\
\altrow
\tcell{Rapport d'incident} & \tcell{Intervention sans interpellation mais avec fait notable} & \tcell{24 heures} \\
\tcell{Rapport d'usage de la force} & \tcell{Tout recours au niveau 3 ou supérieur} & \tcell{24 heures} \\
\altrow
\tcell{Rapport de tir} & \tcell{Tout coup de feu tiré, avec ou sans impact} & \tcell{Immédiat} \\
\tcell{Rapport de poursuite} & \tcell{Toute poursuite véhiculaire} & \tcell{24 heures} \\
\altrow
\tcell{Demande de mandat} & \tcell{Identification d'un suspect en fuite} & \tcell{24 heures} \\
\bottomrule
\end{tabular}

\sstitre{Normes de rédaction}

\begin{itemize}
  \item récit chronologique, au passé, à la première personne ;
  \item faits exclusivement : aucune opinion, aucune supposition, aucun jugement de valeur ;
  \item identité complète et numéro de badge de l'ensemble des agents présents ;
  \item horodatage et localisation précis de chaque étape ;
  \item énoncé du fondement juridique retenu : suspicion raisonnable ou cause probable, et les faits qui la constituent ;
  \item charges alignées sur le Code pénal ;
  \item orthographe et syntaxe soignées : un rapport illisible est un rapport contestable.
\end{itemize}

\sstitre{Mandats}

\cl[mandat-arret]{Une demande de mandat d'arrêt comporte : l'identité du suspect, les charges retenues, le résumé des éléments à charge et l'exposé de la méthode d'identification employée.}

\cl[mandat-perquisition]{Une demande de mandat de perquisition comporte : l'adresse exacte du lieu, l'identité de son occupant, le périmètre exact de la fouille sollicitée, la nature des éléments recherchés et l'exposé de la cause probable qui la justifie.}

\cl{Toute demande de mandat est validée par un membre du Command Staff, puis transmise à l'autorité judiciaire compétente. Un mandat exécuté hors de son périmètre est nul.}

% ---------------------------------------------------------------------
\article{Situations de crise}

\chapeau{Prise d'otages, retranchement, forcené barricadé. Ces situations obéissent à trois principes et à un ordre de priorités qui ne se discutent pas.}

\cl[crise-priorites]{L'ordre des priorités est le suivant, et il est absolu :}

\vspace{4pt}
\noindent
\begin{tabular}{C{16mm} P{132mm}}
\headrow
\thead{Rang} & \thead{Ce qui est protégé} \\
\tcell{\lspdnum{1}} & \tcell{La vie des otages et des tiers présents.} \\
\altrow
\tcell{\lspdnum{2}} & \tcell{La vie des agents engagés.} \\
\tcell{\lspdnum{3}} & \tcell{La vie de l'auteur des faits.} \\
\altrow
\tcell{\lspdnum{4}} & \tcell{L'interpellation de l'auteur.} \\
\tcell{\lspdnum{5}} & \tcell{La préservation des biens.} \\
\bottomrule
\end{tabular}

\cl[crise-principes]{Trois principes gouvernent l'intervention : \stress{contenir}, \stress{temporiser}, \stress{communiquer}. Le temps travaille pour la police. Rien n'impose d'aller vite.}

\sstitre{Répartition des rôles}

\cl[roles]{Dès la caractérisation de la situation, les rôles suivants sont désignés et annoncés :}
\begin{itemize}
  \item \stress{Operation Lead} : superviseur le plus élevé en grade présent. Il commande l'intervention et ne négocie pas lui-même.
  \item \stress{Négociateur} : seul interlocuteur de l'auteur des faits. Il négocie et ne commande pas.
  \item \stress{Négociateur secondaire} : écoute, prend note, alerte le négociateur sur ses angles morts.
\end{itemize}

\cl[validation-nego]{\textbf{Le négociateur ne s'engage sur rien de sa seule autorité.} Toute concession, toute contrepartie et tout accord envisagés avec l'auteur des faits sont soumis à l'Operation Lead et ne prennent effet qu'après sa confirmation expresse. Le négociateur peut temporiser et écouter sans validation ; il ne peut ni promettre ni accorder sans elle.}

\cl[canal-dedie]{Un canal radio dédié est ouvert. Seul le Operation Lead y diffuse. Les unités non engagées maintiennent le silence sur la fréquence principale.}

\sstitre{Conduite de la négociation}

\begin{multicols}{2}
{\lspdsansbk\fontsize{8.6}{10}\selectfont\color{lspdgreen}\addfontfeature{LetterSpace=14}\MakeUppercase{À faire}}\par\vspace{3pt}
\begin{itemize}[leftmargin=4mm,topsep=0pt]
  \item poser des questions ouvertes ;
  \item reformuler ce que dit l'interlocuteur ;
  \item nommer les émotions perçues ;
  \item laisser des silences ;
  \item faire durer, sans jamais mentir.
\end{itemize}
\columnbreak

{\lspdsansbk\fontsize{8.6}{10}\selectfont\color{lspdred}\addfontfeature{LetterSpace=14}\MakeUppercase{À ne pas faire}}\par\vspace{3pt}
\begin{itemize}[leftmargin=4mm,topsep=0pt]
  \item employer le mot « otage » ;
  \item promettre ce qui ne sera pas tenu ;
  \item prendre la colère pour soi ;
  \item minimiser la situation ;
  \item poser un ultimatum sans nécessité.
\end{itemize}
\end{multicols}

\cl[tir-otage]{Lorsqu'un assaut devient inévitable et que des otages sont présents, l'ordre de priorité des cibles est le suivant : l'auteur constituant la menace immédiate, puis le conducteur s'il n'est pas otage, puis les autres auteurs. \textbf{Si la neutralisation d'un auteur risque d'entraîner par réflexe la mort d'un otage, le tir n'est pas engagé.}}

\cl{Les services de secours demeurent en attente hors de la zone dangereuse jusqu'à l'annonce de sécurisation par l'Operation Lead.}

\begin{rappel}[Rythme de l'intervention]
Lent est fluide, fluide est rapide. Une intervention précipitée n'est pas une intervention rapide : c'est une intervention qui échoue. Tant que personne n'est en train de mourir, rien n'impose d'accélérer.
\end{rappel}

% ---------------------------------------------------------------------
\article{Coopération interservices}

\cl[ems]{\stress{Services médicaux} : l'accès des secours est prioritaire en toute circonstance. Les équipes médicales demeurent en attente hors de la zone dangereuse jusqu'à l'annonce de sécurisation par le Operation Lead. La police sécurise, les secours soignent : aucun agent ne s'immisce dans un choix médical. Un blessé placé sous garde le demeure durant sa prise en charge.}

\cl[pompiers]{\stress{Services d'incendie} : sur incendie, explosion ou désincarcération, le commandement de l'action technique appartient aux services d'incendie. La police établit le périmètre, régule la circulation et protège l'intervention.}

\cl[doj]{\stress{Autorité judiciaire} : le ministère public est l'interlocuteur exclusif pour la validation des mandats, les réquisitions et les négociations de peine. \textbf{Aucun agent ne négocie directement une peine avec un mis en cause ou son conseil.}}

\cl[avocats]{\stress{Avocats} : l'accès au client placé en garde à vue est garanti, de même que la confidentialité de leur entretien. La communication des pièces s'effectue par la voie institutionnelle.}

\cl[agences]{\stress{Autres agences} : la compétence territoriale appartient à l'agence du lieu. L'agence ayant initié une poursuite ou une procédure en conserve la direction jusqu'à passation formelle annoncée à la radio.}

\cl{\stress{Dépanneurs et fourrière} : l'agent demeure sur place jusqu'à l'enlèvement du véhicule, relève l'immatriculation et consigne le motif de l'enlèvement.}

\cl[conflit-competence]{Tout conflit de compétence se règle entre superviseurs, hors de la fréquence ouverte et hors de la présence du public.}

% ---------------------------------------------------------------------
\article{Accompagnement civil en patrouille}

\cl[ridealong]{Les journalistes, personnels des services de secours, magistrats et auxiliaires de justice peuvent être autorisés à accompagner une patrouille sur autorisation du Command Staff. Cette autorisation fait l'objet d'une note de service établissant :}
\begin{enumerate}
  \item qu'il peut être objectivement attesté que le tiers ne mettra en danger ni la vie des citoyens, ni celle des agents qu'il accompagne ;
  \item que le tiers est physiquement et psychologiquement en mesure de faire face à un événement violent ;
  \item la liste des règles que le tiers s'engage à respecter et dont il atteste avoir pris connaissance avant le début de la patrouille.
\end{enumerate}

\cl[casier-ridealong]{Toute personne ne présentant pas un casier vierge est automatiquement refusée. Une vérification des antécédents et de l'absence de mandat en cours est effectuée préalablement.}

\cl{L'accompagnant ne porte aucune arme. Le port d'un gilet de protection fourni par le département est obligatoire. Il n'accède ni à la radio, ni au terminal embarqué.}

\cl[ridealong-danger]{En cas d'intervention violente, de poursuite ou d'appel général, l'accompagnant est déposé au poste le plus proche ou en lieu sûr avant tout engagement de l'unité.}

\cl{Aucun agent ne peut prendre en accompagnement un membre de sa famille ou une personne avec laquelle il entretient une relation personnelle.}

\cl{Aucun enregistrement, image ou compte rendu établi lors d'un accompagnement ne peut être diffusé sans l'accord préalable du Command Staff.}

\begin{interdit}[Détournement de l'accompagnement]
L'accompagnement en patrouille ne constitue en aucun cas un mode d'accès aux fonctions de police, ni un substitut à la formation dispensée par l'académie. Il ne confère aucune prérogative, aucun statut et aucune priorité de recrutement.
\end{interdit}
